\chapter{Volatilité stochastique}\label{ch:dv:stochastic-volatility}

Interrogée sur la raison pour laquelle la corrélation calibrée de son modèle sur indice vaut $-0.7$, une trader d'options répond
sans formule~: parce que lorsque l'indice baisse, la volatilité monte. Interrogée sur la raison pour laquelle la \omterm{def:dv:stochastic-volatility:sv}{volatilité de
la volatilité} vaut 0.6, elle désigne les ailes du smile à un an, que le modèle ne peut atteindre que si la variance elle-même fluctue fortement.
Interrogée sur la raison pour laquelle la vitesse de retour à la moyenne vaut 2, elle trace la
structure par terme à la monnaie qui revient vers son niveau de long terme en moins d'un an. Chacun des
cinq nombres du \omterm{def:dv:stochastic-volatility:heston}{modèle de Heston} a un sens sur le desk, et c'est pour cela qu'il a survécu
trente ans d'alternatives mieux ajustées. Ce chapitre introduit la volatilité stochastique,
dérive ce qui peut l'être du \omterm{def:dv:stochastic-volatility:heston}{modèle de Heston}, évalue les options à l'aide de sa
fonction caractéristique, calibre le modèle sur la \omterm{def:dv:implied-volatility-and-its-surface:surface}{surface de volatilité} du chapitre~9 et lit ses paramètres comme le ferait un trader.
Il mesure aussi où il échoue~: aux échéances courtes, son skew n'est que la moitié de
ce que montre le marché.

\section{Modèles à volatilité stochastique}

La \omterm{def:dv:local-volatility:model}{volatilité locale} fait de la volatilité une fonction du spot, et relie donc chaque mouvement du smile à
un mouvement du spot. La volatilité réelle a une vie propre~: elle monte après les baisses, mais elle
monte et baisse aussi les jours où le spot ne bouge pas. L'effet de levier (le livre~4,
chapitre~18) est le premier fait, mesuré comme une corrélation négative entre les rendements et la
volatilité ultérieure, forte et brève pour les indices et plus faible pour les actions individuelles~; le regroupement de volatilité
est le second.

\begin{definition}[Modèle à volatilité stochastique]\label{def:dv:stochastic-volatility:sv}
Un \emph{modèle à volatilité stochastique}\index{modèle à volatilité stochastique} précise, sous la
mesure de valorisation,
\[
dS_t=(r-q)S_t\,dt+\sqrt{v_t}\,S_t\,dW^1_t,\qquad dv_t=\alpha(v_t)\,dt+\beta(v_t)\,dW^2_t,\qquad
d\langle W^1,W^2\rangle_t=\rho\,dt,
\]
où la variance instantanée $v_t$ est une seconde variable d'état. La
\emph{corrélation spot-volatilité}\index{corrélation spot-volatilité} $\rho$ relie les deux
chocs, et l'échelle de $\beta$ par rapport à $v$ est la \emph{volatilité de la
volatilité}\index{volatilité de la volatilité}.
\end{definition}

Deux sources d'aléa et un actif risqué négocié rendent le marché incomplet (chapitre~1)~: une
option ne peut être répliquée avec l'action et le cash seuls. Elle peut l'être avec l'action et une seconde
option, qui permet de couvrir le risque de variance. La dérive $\alpha$ sous la mesure de valorisation absorbe un
prix de marché du risque de volatilité que les dynamiques réelles ne déterminent pas~; la calibration sur
les prix d'options, et non une régression sur l'historique, la fixe.

\section{Le modèle de Heston}

\begin{definition}[Modèle de Heston]\label{def:dv:stochastic-volatility:heston}
Le \emph{modèle de Heston}\index{modèle de Heston} suppose que la variance suit un processus en racine carrée (le
livre~4, chapitre~4)~:
\[
dv_t=\kappa(\bar v-v_t)\,dt+\eta\sqrt{v_t}\,dW^2_t,
\]
avec cinq paramètres~: la variance initiale $v_0$, la vitesse de retour à la moyenne $\kappa$, la variance de long
terme $\bar v$, la \omterm{def:dv:stochastic-volatility:sv}{volatilité de la volatilité} $\eta$ et la corrélation $\rho$. La variance reste
non négative et n'atteint jamais zéro lorsque la condition de Feller $2\kappa\bar v\ge\eta^2$ est vérifiée.
\end{definition}

La valeur du modèle tient à son caractère affine~: le log-prix et la variance entrent dans l'équation de valorisation
linéairement dans l'exposant, et sa fonction caractéristique est connue sous forme fermée.

\begin{proposition}[Fonction caractéristique de Heston]\label{prop:dv:stochastic-volatility:cf}
Soit $x_T=\ln(S_T/F_{0,T})$. Sa fonction caractéristique (le livre~4, chapitre~1) est
$\varphi_{x_T}(u)=\exp\bigl(C(u)+D(u)v_0\bigr)$, avec $\xi=\kappa-\rho\eta iu$,
$d=\sqrt{\xi^2+\eta^2(iu+u^2)}$, $g=(\xi-d)/(\xi+d)$ et
\[
C=\frac{\kappa\bar v}{\eta^2}\Bigl((\xi-d)T-2\ln\frac{1-ge^{-dT}}{1-g}\Bigr),\qquad
D=\frac{\xi-d}{\eta^2}\,\frac{1-e^{-dT}}{1-ge^{-dT}} .
\]
\end{proposition}
\admitted

Les fonctions $C$ et $D$ résolvent deux équations différentielles ordinaires (une équation de Riccati et une
intégrale de celle-ci) obtenues en substituant une forme exponentielle-affine supposée dans l'équation de valorisation~;
Heston les a résolues en 1993. La forme ci-dessus conserve le logarithme complexe sur sa branche principale pour
tout $u$~; la forme originale de Heston, algébriquement égale, saute d'une branche à l'autre pour les échéances longues et
donne des prix erronés, un défaut connu sous le nom de «~little Heston trap~».

\section{Fonction caractéristique et valorisation}

Grâce à la fonction caractéristique, une option européenne se réduit à une seule intégrale. En termes de forward, avec
$k=\ln(K/F)$,
\begin{equation}\label{eq:dv:stochastic-volatility:lewis}
C(K,T)=P(0,T)\Bigl(F-\frac{\sqrt{FK}}{\pi}\int_0^\infty\operatorname{Re}\bigl(e^{-iuk}
\varphi_{x_T}(u-\tfrac i2)\bigr)\frac{du}{u^2+\frac14}\Bigr),
\end{equation}
une formule que le chapitre~24 établit et généralise à tout modèle dont la fonction caractéristique est connue.
Tous les strikes d'une même échéance partagent l'évaluation de $\varphi$, si bien qu'un smile entier coûte une
intégrale vectorisée. Trois contrôles vérifient l'implémentation~: $\varphi(0)=\varphi(-i)=1$ (le
forward est une martingale)~; lorsque $\eta\to0$ le smile s'aplatit à la racine de la variance moyenne de la
trajectoire déterministe~; une simulation des deux équations reproduit les prix. Avec les paramètres de base du
chapitre ($v_0=\bar v=0.04$, $\kappa=1.5$, $\eta=0.6$, $\rho=-0.7$), le smile à un an va de
23.4~\% au strike 80 à 13.3~\% au 120, et une simulation de 200\,000 trajectoires concorde avec l'intégrale dans les
limites de son bruit statistique.

\section{Calibration}

Calibrer Heston sur une surface est un problème de moindres carrés non linéaires à cinq paramètres, résolu par
l'algorithme de Levenberg--Marquardt (le livre~4, chapitre~24) sur les résidus de volatilité implicite, avec
paramètres transformés (logarithmes, tangente hyperbolique pour $\rho$) de sorte qu'aucun pas ne fasse
sortir de l'ensemble admissible.

\begin{example}[Heston sur la surface du chapitre~9]\label{ex:dv:stochastic-volatility:calib}
Calibré sur neuf strikes (80 à 120) à quatre échéances (trois mois à deux ans) de la surface de
chapitre~9, Heston donne $v_0=0.0267$ (une volatilité spot de 16.3~\%), $\kappa=2.05$, $\bar v=0.060$
(24.5~\% de long terme), $\eta=0.66$ et $\rho=-0.70$, avec une erreur quadratique moyenne de 0.37 point de volatilité
et au plus 0.7 point, à deux ans (voir la \cref{fig:dv:stochastic-volatility:calib}). La condition
de Feller n'est pas vérifiée, $2\kappa\bar v-\eta^2=-0.19$~: la variance ajustée peut toucher zéro, comme dans la plupart des calibrations sur
indices.
\end{example}

\begin{omfigure}
\begin{tikzpicture}
\begin{axis}[width=10.5cm,height=5cm,
  xlabel={strike},ylabel={volatilité implicite (\%)},
  tick label style={font=\small},label style={font=\small},
  xmin=80,xmax=120,ymin=12,ymax=28,grid=major,grid style={gray!25},
  legend columns=2,legend style={font=\footnotesize,at={(0.5,-0.3)},anchor=north,draw=none,fill=none,/tikz/every even column/.append style={column sep=8pt}},
  table/col sep=comma]
\addplot[omDef,very thick] table[x=strike,y=mkt3m]{figdata/derivatives/10-stochastic-volatility/calib.csv};
\addplot[omDef,thick,dashed] table[x=strike,y=mod3m]{figdata/derivatives/10-stochastic-volatility/calib.csv};
\addplot[omThm,very thick] table[x=strike,y=mkt2y]{figdata/derivatives/10-stochastic-volatility/calib.csv};
\addplot[omThm,thick,dashed] table[x=strike,y=mod2y]{figdata/derivatives/10-stochastic-volatility/calib.csv};
\legend{marché 3 mois,Heston 3 mois,marché 2 ans,Heston 2 ans}
\end{axis}
\end{tikzpicture}

\omcaption{Le \omterm{def:dv:stochastic-volatility:heston}{modèle de Heston} calibré (en pointillés) face à la surface à laquelle il a été ajusté (en trait plein) à trois
mois et deux ans. Cinq paramètres ajustent quatre échéances à 0.7 point de volatilité près. Données~: le
tutoriel.}\label{fig:dv:stochastic-volatility:calib}
\end{omfigure}

L'ajustement masque un écart structurel qui apparaît aux échéances hors de l'ensemble de calibration. Le skew à la monnaie de Heston
tend vers une limite finie lorsque l'échéance se réduit, tandis que les skews de marché continuent de se raidir, à peu près comme
$T^{-1/2}$ dans cette surface. À une semaine, le skew du modèle calibré vaut la moitié de celui du marché
(\cref{fig:dv:stochastic-volatility:skew})~; les sauts (chapitre~13) et la volatilité rugueuse (chapitre~12) sont les
deux réponses.

\begin{omfigure}
\begin{tikzpicture}
\begin{loglogaxis}[width=10.5cm,height=5cm,
  xlabel={échéance (en années, échelle logarithmique)},ylabel={$-\partial\sigma_{\mathrm{imp}}/\partial k$ à la monnaie},
  tick label style={font=\small},label style={font=\small},
  xmin=0.015,xmax=2.5,ymin=0.12,ymax=2,grid=major,grid style={gray!25},
  ytick={0.15,0.2,0.3,0.5,0.7,1,1.5},yticklabels={0.15,0.2,0.3,0.5,0.7,1,1.5},
  xtick={0.02,0.05,0.1,0.25,0.5,1,2},xticklabels={0.02,0.05,0.1,0.25,0.5,1,2},
  legend columns=2,legend style={font=\footnotesize,at={(0.5,-0.3)},anchor=north,draw=none,fill=none,/tikz/every even column/.append style={column sep=8pt}},
  table/col sep=comma]
\addplot[omThm,very thick,mark=*,mark size=1.5pt] table[x=t,y=market]{figdata/derivatives/10-stochastic-volatility/skew_term.csv};
\addplot[omDef,very thick,mark=square*,mark size=1.5pt] table[x=t,y=heston]{figdata/derivatives/10-stochastic-volatility/skew_term.csv};
\legend{surface de marché,Heston calibré}
\end{loglogaxis}
\end{tikzpicture}

\omcaption{Skew à la monnaie par échéance, en échelles logarithmiques~: le skew de la surface de marché continue de se raidir aux échéances
courtes, le skew du \omterm{def:dv:stochastic-volatility:heston}{modèle de Heston} calibré se stabilise autour de 0.7. À une semaine, celui du modèle vaut la moitié de celui du
marché~; entre trois mois et un an, ils se confondent, c'est-à-dire là où le modèle a été calibré. Données~: le
tutoriel.}\label{fig:dv:stochastic-volatility:skew}
\end{omfigure}

\section{Ce que les paramètres signifient pour un trader}

Chaque paramètre déplace une caractéristique de la surface
(\cref{fig:dv:stochastic-volatility:params,fig:dv:stochastic-volatility:term})~:

\begin{itemize}
\item $\rho$ \emph{incline} le smile~: plus négatif, skew plus raide. Faire passer $\rho$ de $-0.7$ à $-0.5$
abaisse le skew à un an 90--110 de 5.71 à 3.94 points de volatilité.
\item $\eta$ \emph{courbe} le smile~: les ailes montent avec la \omterm{def:dv:stochastic-volatility:sv}{volatilité de la volatilité}. La courbure à un an
mesure $\tfrac12(\sigma_{90}+\sigma_{110})-\sigma_{100}$ vaut 0.30 point pour $\eta=0.6$~; une courbure de 0.5
point exige $\eta=0.77$.
\item $v_0$ et $\bar v$ fixent les extrémités courtes et longues de la structure par terme à la monnaie, et $\kappa$ la
vitesse à laquelle l'une devient l'autre~; $\kappa$ fait aussi décroître le skew avec l'échéance, puisque la variance
est soumise à des chocs qui s'oublient à ce rythme.
\end{itemize}

\begin{omfigure}
\begin{tikzpicture}
\begin{groupplot}[group style={group size=2 by 1,horizontal sep=1.5cm},
  width=6.6cm,height=4.8cm,xlabel={strike},
  tick label style={font=\small},label style={font=\small},
  xmin=70,xmax=130,ymin=8,ymax=30,grid=major,grid style={gray!25},table/col sep=comma]
\nextgroupplot[ylabel={volatilité implicite à un an (\%)},title={\small corrélation $\rho$},
  legend columns=3,legend style={font=\footnotesize,at={(0.5,-0.3)},anchor=north,draw=none,fill=none}]
\addplot[omThm,very thick] table[x=strike,y=rho_m09]{figdata/derivatives/10-stochastic-volatility/params.csv};
\addplot[omDef,thick] table[x=strike,y=rho_m05]{figdata/derivatives/10-stochastic-volatility/params.csv};
\addplot[omProp,thick,dashed] table[x=strike,y=rho_0]{figdata/derivatives/10-stochastic-volatility/params.csv};
\legend{$-0.9$,$-0.5$,$0$}
\nextgroupplot[title={\small volatilité de la volatilité $\eta$},
  legend columns=3,legend style={font=\footnotesize,at={(0.5,-0.3)},anchor=north,draw=none,fill=none}]
\addplot[omProp,thick,dashed] table[x=strike,y=eta03]{figdata/derivatives/10-stochastic-volatility/params.csv};
\addplot[omDef,thick] table[x=strike,y=eta06]{figdata/derivatives/10-stochastic-volatility/params.csv};
\addplot[omThm,very thick] table[x=strike,y=eta09]{figdata/derivatives/10-stochastic-volatility/params.csv};
\legend{$0.3$,$0.6$,$0.9$}
\end{groupplot}
\end{tikzpicture}

\omcaption{Smiles de Heston à un an autour des paramètres de base ($v_0=\bar v=0.04$, $\kappa=1.5$,
$\eta=0.6$, $\rho=-0.7$)~: la corrélation fait pivoter le smile (à gauche), la \omterm{def:dv:stochastic-volatility:sv}{volatilité de la volatilité} relève
ses ailes et, avec une corrélation négative, son skew (à droite). Données~: le code du
chapitre.}\label{fig:dv:stochastic-volatility:params}
\end{omfigure}

\begin{omfigure}
\begin{tikzpicture}
\begin{axis}[width=10.5cm,height=4.8cm,
  xlabel={échéance (en années)},ylabel={volatilité à la monnaie (\%)},
  tick label style={font=\small},label style={font=\small},
  xmin=0,xmax=5,ymin=12,ymax=29,grid=major,grid style={gray!25},
  legend columns=3,legend style={font=\footnotesize,at={(0.5,-0.3)},anchor=north,draw=none,fill=none,/tikz/every even column/.append style={column sep=8pt}},
  table/col sep=comma]
\addplot[omDef,thick,mark=*,mark size=1.3pt] table[x=t,y=v0_002]{figdata/derivatives/10-stochastic-volatility/term.csv};
\addplot[gray,thick,mark=square*,mark size=1.3pt] table[x=t,y=v0_004]{figdata/derivatives/10-stochastic-volatility/term.csv};
\addplot[omThm,thick,mark=triangle*,mark size=1.6pt] table[x=t,y=v0_008]{figdata/derivatives/10-stochastic-volatility/term.csv};
\legend{$v_0=0.02$,$v_0=0.04$,$v_0=0.08$}
\end{axis}
\end{tikzpicture}

\omcaption{Structures par terme à la monnaie pour trois variances initiales, la variance de long terme valant 0.04
(20~\%) et $\kappa=1.5$~: la structure par terme croissante d'un marché calme et la structure inversée d'un marché tendu,
toutes deux convergeant à la vitesse $\kappa$. Avec $v_0=\bar v$, la courbe n'est pas plate~: elle descend de 19.6~\% à 16.9~\% à
un an, car le prix à la monnaie est presque linéaire en volatilité et non en variance, et le moyennage sur une
variance aléatoire abaisse la volatilité moyenne (la \omterm{prop:dv:stochastic-volatility:mixing}{formule de mélange} ci-dessous)~; la corrélation négative l'abaisse davantage. Données~:
le code du chapitre.}\label{fig:dv:stochastic-volatility:term}
\end{omfigure}

Pour un trader, l'identité la plus utile est celle qui tient sans corrélation.

\begin{proposition}[Formule de mélange]\label{prop:dv:stochastic-volatility:mixing}
Si $\rho=0$, une option européenne dans un \omterm{def:dv:stochastic-volatility:sv}{modèle à volatilité stochastique} vaut le prix de Black--Scholes évalué
à la racine de la variance moyenne de la trajectoire, moyenné sur les trajectoires de variance~:
$C=\E\bigl[C_{\mathrm{BS}}(\bar\sigma)\bigr]$ avec $\bar\sigma^2=\frac1T\int_0^Tv_t\,dt$. Pour $\rho\ne0$, la même
relation reste valable en remplaçant le spot par un spot ajusté et la variance par $(1-\rho^2)$ fois celle-ci~; c'est la
\emph{formule de mélange}\index{formule de mélange}.
\end{proposition}
\begin{proof}[Démonstration pour $\rho=0$]
Conditionnellement à toute la trajectoire de variance, $\ln S_T$ suit une loi normale de variance $\int_0^Tv_t\,dt$, car
$W^1$ est indépendant de $W^2$~: le prix conditionnel est celui de Black--Scholes avec cette variance. Il suffit de prendre
l'espérance.
\end{proof}

La formule explique d'emblée la forme du smile~: les prix de Black--Scholes sont convexes en volatilité hors de
la monnaie, si bien que la moyenne sur une volatilité aléatoire relève davantage les prix hors de la monnaie que ceux à la monnaie,
et le smile se courbe vers le haut (la \omterm{def:dv:greeks-and-the-hedging-pnl:greeks}{volga} du chapitre~4 à l'œuvre). Une simulation de trajectoires de variance avec
$\rho=0$ donne, par la formule, 7.136 pour le call à la monnaie d'un an contre 7.137 par la
fonction caractéristique, à l'intérieur de son erreur type de 0.011.

\section{Tutoriel~: pricing et calibration de Heston}\label{tut:dv:stochastic-volatility:tut}

\begin{tutorial}
\textbf{Objectif.} Pricer avec la fonction caractéristique, la vérifier de trois façons, calibrer sur une surface et lire
les paramètres. \textbf{État final~:} les quatre figures du chapitre et les nombres de
l'\cref{ex:dv:stochastic-volatility:calib}.
\begin{tutsteps}
\item \textbf{La fonction caractéristique et l'intégrale}, un smile à la fois~:
\omcode{code/firm/heston/firm_heston.py}{32}{54}{Fonction caractéristique de Heston (forme à branche principale) et prix de Fourier d'un smile.}\label{lst:dv:stochastic-volatility:cf}
\item \textbf{Levenberg--Marquardt} sur les résidus de volatilité implicite, avec des paramètres transformés et un
amortissement qui s'adapte au succès~:
\omcode{code/firm/heston/firm_heston.py}{74}{104}{Calibration des cinq paramètres.}\label{lst:dv:stochastic-volatility:calib}
\item \textbf{Exécuter} \texttt{dv\_heston.calibration()}, \texttt{mixing\_check()} et \texttt{fig\_heston.py}.
\end{tutsteps}
\textbf{Pour aller plus loin.} Remplacer la forme à branche principale par la forme originale de Heston et pricer une
option à cinq ans avec une grande \omterm{def:dv:stochastic-volatility:sv}{volatilité de la volatilité}~: repérer les strikes où le prix saute~; ajouter alors les échéances à une semaine et
à un mois à la calibration et observer la dégradation de l'ajustement ailleurs.
\end{tutorial}

\section{Construction~: le modèle de Heston}\label{bld:dv:stochastic-volatility:heston}

\begin{build}
\bfield{Rôle} Premier modèle du mini-cabinet doté de sa propre dynamique de volatilité~: il price des produits qui
dépendent des \omterm{def:dv:local-volatility:forward}{smiles forward} (chapitre~16), il constitue la partie stochastique de la volatilité stochastique-locale (chapitre~20),
et il sert de cas de test aux moteurs de Fourier et de calibration (chapitre~24).
\bfield{Interface} \texttt{Heston(v0, kappa, vbar, eta, rho)} avec \texttt{cf(u, t)} et \texttt{feller()}~;
\texttt{call\_prices(model, fwd, strikes, t, df)}~; \texttt{implied\_vols}~; \texttt{calibrate(quotes, fwd,
start) -> (model, rmse)}~; \texttt{simulate(model, s0, t, steps, n\_paths, seed)}.
\bfield{Règles} Fonction caractéristique sous la forme à branche principale~; le forward comme numéraire~; calibration sur
les résidus de volatilité implicite avec des paramètres transformés~; signaler l'écart de Feller, ne jamais l'imposer en silence.
\bfield{Tests d'acceptation} \texttt{code/firm/heston/tests/}~: $\varphi(0)=\varphi(-i)=1$~; la limite en $\eta$ petit
égale Black--Scholes à la variance moyenne déterministe~; la simulation concorde avec l'intégrale à moins de
trois erreurs types~; la calibration retrouve des paramètres connus à partir de leurs propres smiles.
\bfield{Pour aller plus loin} Le schéma de simulation quadratique-exponentiel (chapitre~23) et les gradients analytiques des prix
par rapport aux paramètres, pour une calibration plus rapide.
\end{build}

\begin{omsources}
\item S.~L.~Heston, ``A closed-form solution for options with stochastic volatility with applications to bond
and currency options'', \emph{Review of Financial Studies} 6 (1993) 327--343.
\item H.~Albrecher, P.~Mayer, W.~Schoutens and J.~Tistaert, ``The little Heston trap'', \emph{Wilmott} (2007)
83--92.
\item J.~Hull and A.~White, ``The pricing of options on assets with stochastic volatilities'', \emph{Journal of
Finance} 42 (1987) 281--300.
\item M.~Romano and N.~Touzi, ``Contingent claims and market completeness in a stochastic volatility model'',
\emph{Mathematical Finance} 7 (1997) 399--412.
\item J.-P.~Bouchaud, A.~Matacz and M.~Potters, ``Leverage effect in financial markets: the retarded
volatility model'', \emph{Physical Review Letters} 87 (2001) 228701.
\item J.~Gatheral, \emph{The Volatility Surface}, Wiley, 2006, chapitres 2--3.
\end{omsources}

\section{Exercices}

\begin{exercise}[$\star$]\label{exo:dv:stochastic-volatility:1}
Vérifier la condition de Feller pour les paramètres de base et pour les paramètres calibrés.
\end{exercise}

\begin{exercise}[$\star$]\label{exo:dv:stochastic-volatility:2}
Avec $v_0=0.08$, $\bar v=0.04$ et $\kappa=1.5$, donner la variance déterministe après six mois et après un an,
ainsi que la volatilité à la monnaie d'une option à un an lorsque $\eta\to0$.
\end{exercise}

\begin{exercise}[$\star$]\label{exo:dv:stochastic-volatility:3}
Pourquoi doit-on avoir $\varphi_{x_T}(-i)=1$~? Que signifierait une valeur de 1.01~?
\end{exercise}

\begin{exercise}[$\star\star$]\label{exo:dv:stochastic-volatility:4}
Pourquoi le prix de marché du risque de volatilité n'apparaît-il pas dans le modèle calibré, et qu'implique cela
pour l'usage de la volatilité historique afin de fixer $\bar v$~?
\end{exercise}

\begin{exercise}[$\star\star$]\label{exo:dv:stochastic-volatility:5}
Utiliser la \omterm{prop:dv:stochastic-volatility:mixing}{formule de mélange} pour expliquer pourquoi, avec $\rho=0$, le smile de Heston est symétrique en \omterm{def:dv:implied-volatility-and-its-surface:surface}{log-moneyness} autour du
forward et pourquoi sa volatilité à la monnaie se situe sous la variance moyenne.
\end{exercise}

\begin{exercise}[$\star\star$]\label{exo:dv:stochastic-volatility:6}
Quel paramètre faudrait-il déplacer pour ajuster un marché où le skew à un mois s'est raidi du jour au lendemain alors que le skew à un an
n'a pas bougé~? Pourquoi Heston ne peut-il pas l'ajuster par un seul déplacement de paramètre~?
\end{exercise}

\begin{exercise}[$\star\star\star$]\label{exo:dv:stochastic-volatility:7}
\emph{Programmation.} Avec \texttt{eta\_for\_butterfly}, trouver la \omterm{def:dv:stochastic-volatility:sv}{volatilité de la volatilité} qui donne une courbure à un an
de 0.5 point de volatilité, ainsi que le skew 90--110 à un an pour cette valeur de $\eta$.
\end{exercise}

\begin{exercise}[$\star\star\star$]\label{exo:dv:stochastic-volatility:8}
\emph{Trouver l'erreur.} «~Notre calibration de Heston ajuste la surface de trois mois à deux ans à 0.4 point près, nous
l'utilisons donc pour pricer les options hebdomadaires sur l'indice.~»
\end{exercise}

\section{Problème~: lecture des paramètres}

\begin{problem}[{Problème du week-end --- ce que chaque nombre de Heston fait au smile}]\label{pb:dv:stochastic-volatility:1}
Un trader junior hérite du \omterm{def:dv:stochastic-volatility:heston}{modèle de Heston} de base ($v_0=\bar v=0.04$, $\kappa=1.5$, $\eta=0.6$, $\rho=-0.7$, taux nuls)
et doit l'ajuster à deux changements du marché~: le skew 90--110 à un an baisse d'environ 1.8 point, et la
courbure à un an monte à 0.5 point.

\medskip\noindent\textbf{Partie I --- Le smile de base.}
\begin{enumerate}
\item Donner les volatilités implicites à un an pour les strikes 80, 90, 100, 110 et 120.
\item Donner le skew 90--110 à un an ainsi que la mesure de courbure.
\item La condition de Feller est-elle satisfaite~?
\item Pourquoi la volatilité à la monnaie est-elle inférieure à 20~\% alors que $v_0=\bar v=0.04$~?
\item Que prédit le modèle pour la volatilité à la monnaie à cinq ans~?
\end{enumerate}

\medskip\noindent\textbf{Partie II --- La corrélation.}
\begin{enumerate}[resume]
\item Donner le skew 90--110 avec $\rho=-0.5$.
\item Donner la variation du skew pour une variation de 0.2 de la corrélation.
\item La volatilité à la monnaie varie-t-elle lorsque $\rho$ change~? Pourquoi~?
\item Quels produits sont les plus exposés à $\rho$~?
\item Que signifie $\rho=-0.7$ pour l'indice lorsqu'il chute de 1~\% en une journée~?
\end{enumerate}

\medskip\noindent\textbf{Partie III --- La \omterm{def:dv:stochastic-volatility:sv}{volatilité de la volatilité}.}
\begin{enumerate}[resume]
\item Donner le $\eta$ qui produit une courbure à un an de 0.5 point.
\item Que devient le skew pour ce $\eta$, et pourquoi~?
\item Comment restaureriez-vous le skew~?
\item Quel est l'écart de Feller pour ce $\eta$~?
\item Quels produits sont les plus exposés à $\eta$~?
\end{enumerate}

\medskip\noindent\textbf{Partie IV --- Jugement.}
\begin{enumerate}[resume]
\item Pourquoi le trader devrait-il recalibrer les cinq paramètres plutôt que d'en modifier un à la fois~?
\item Pourquoi le desk pourrait-il garder $\kappa$ fixe d'un jour à l'autre~?
\item Que le modèle fait-il de faux à une semaine~?
\item Énoncer le \emph{résultat nommé}~: la variation du skew 90--110 à un an pour une variation de $+0.2$ de la corrélation,
et la \omterm{def:dv:stochastic-volatility:sv}{volatilité de la volatilité} qui correspond à une courbure de 0.5 point.
\item En une phrase~: que valorise la \omterm{def:dv:stochastic-volatility:sv}{volatilité de la volatilité}~?
\end{enumerate}
\end{problem}

\section{Questions d'entretien}

\begin{interviewq}[$\star$ \iqroles{trader, chercheur}]\label{iq:dv:stochastic-volatility:1}
Écrivez le \omterm{def:dv:stochastic-volatility:heston}{modèle de Heston} et dites ce que fait chaque paramètre au smile.
\end{interviewq}

\begin{interviewq}[$\star$ \iqroles{chercheur}]\label{iq:dv:stochastic-volatility:2}
Pourquoi un marché à volatilité stochastique est-il incomplet, et comment couvririez-vous une option dans un tel marché~?
\end{interviewq}

\begin{interviewq}[$\star\star$ \iqroles{chercheur}]\label{iq:dv:stochastic-volatility:3}
Établissez le prix d'une option européenne dans un \omterm{def:dv:stochastic-volatility:sv}{modèle à volatilité stochastique} de corrélation nulle.
\end{interviewq}

\begin{interviewq}[$\star\star$ \iqroles{développeur, chercheur}]\label{iq:dv:stochastic-volatility:4}
Votre pricer de Heston renvoie des prix faux pour les longues échéances avec une grande \omterm{def:dv:stochastic-volatility:sv}{volatilité de la volatilité}. Que soupçonnez-vous~?
\end{interviewq}

\begin{interviewq}[$\star\star$ \iqroles{trader, risque}]\label{iq:dv:stochastic-volatility:5}
La corrélation de Heston calibrée est passée de $-0.7$ à $-0.4$ pendant la nuit, avec la même erreur d'ajustement. Que faites-vous~?
\end{interviewq}

\begin{interviewq}[$\star\star\star$ \iqroles{chercheur}]\label{iq:dv:stochastic-volatility:6}
Pourquoi le skew de Heston est-il trop plat aux courtes échéances~? Citez deux façons de le corriger.
\end{interviewq}
